\documentclass[14pt, a4paper]{extarticle}
\usepackage[utf8]{inputenc}
\usepackage{cmap}
\usepackage[T2A]{fontenc}
\usepackage[russian]{babel}
\usepackage{amsmath}
\usepackage{amssymb}
\usepackage{amsthm}
\usepackage{geometry}
\usepackage{breqn}
\allowdisplaybreaks
\geometry{
    top=2cm,
    bottom=2cm,
    left=3cm,
    right=2cm
}

\begin{document}

\title{Математическое обоснование работы патрейсера на видеокарте}
\author{}
\date{}
\maketitle

\section{Введение и архитектура}

\subsection{Обзор GPU патрейсинга}

Path tracing (патрейсинг) --- это метод фотореалистичного рендеринга, основанный на решении уравнения рендеринга методом Монте-Карло. Выполнение патрейсинга на GPU (Graphics Processing Unit) позволяет значительно ускорить вычисления благодаря массивному параллелизму.

Основная идея заключается в том, что каждый пиксель изображения обрабатывается независимо, что идеально подходит для параллельной архитектуры GPU. Каждый поток GPU вычисляет цвет одного пикселя, трассируя множество случайных лучей и усредняя результаты.

\subsection{Архитектура OpenGL Compute Shaders}

В данной реализации используется OpenGL 4.3+ с поддержкой compute shaders. Compute shader --- это специальный тип шейдера, который позволяет выполнять произвольные вычисления на GPU без привязки к графическому конвейеру.

Архитектура compute shader основана на концепции work groups (рабочих групп). Каждая рабочая группа содержит фиксированное количество потоков (threads), которые выполняются параллельно. В нашем случае используется размер рабочей группы $16 \times 16 = 256$ потоков.

\subsection{Структура данных на GPU}

Данные передаются на GPU через специальные буферы:
\begin{itemize}
    \item \textbf{Uniform Buffer Objects (UBO)} --- для постоянных данных (камера, небо, туман)
    \item \textbf{Shader Storage Buffer Objects (SSBO)} --- для массивов данных (материалы, объекты, накопительный буфер)
    \item \textbf{Image2D} --- для выходного изображения (текстура RGBA16F)
\end{itemize}

\section{Форматы данных и передача на GPU}

\subsection{Uniform Buffer Objects (UBO) --- формат std140}

Uniform Buffer Objects используют выравнивание std140, которое гарантирует определенную структуру размещения данных в памяти. Это необходимо для корректного доступа к данным из шейдера.

\subsubsection{Блок камеры}

Камера передается как структура из 16 float32 (4 vec4), что составляет 64 байта:

\begin{align}
\text{CameraBlock} &= \begin{bmatrix}
\text{camPos} & \text{(vec4, 16 байт)} \\
\text{camTarget} & \text{(vec4, 16 байт)} \\
\text{camUp} & \text{(vec4, 16 байт)} \\
\text{camFov, camAperture, camFocusDist, camAspect} & \text{(4 float, 16 байт)}
\end{bmatrix}
\end{align}

Математически позиция камеры $\mathbf{c}_{\text{pos}}$, цель $\mathbf{c}_{\text{target}}$ и вектор вверх $\mathbf{c}_{\text{up}}$ определяют систему координат камеры.

\subsubsection{Блок неба}

Небо передается как структура с выравниванием std140:

\begin{align}
\text{SkyBlock} &= \begin{bmatrix}
\text{skyType} & \text{(float, 4 байта + 12 байт padding)} \\
\text{skyColor} & \text{(vec4, 16 байт)} \\
\text{skyHorizon} & \text{(vec4, 16 байт)} \\
\text{skyZenith} & \text{(vec4, 16 байт)}
\end{bmatrix}
\end{align}

Общий размер: 64 байта (выравнивание std140 требует, чтобы каждый элемент занимал кратные 16 байт).

\subsubsection{Блок тумана}

Туман содержит физические параметры объемного рассеяния:

\begin{align}
\text{FogBlock} &= \begin{bmatrix}
\text{fogDensity, fogScatter, fogAffectSky, fogGpuVolumetric} & \text{(4 float)} \\
\text{fogColor} & \text{(vec4)} \\
\text{fogSigmaS, fogSigmaA, fogG} & \text{(3 float + padding)} \\
\text{fogHeteroStrength, fogNoiseScale, fogNoiseOctaves} & \text{(3 float + padding)}
\end{bmatrix}
\end{align}

Ключевые физические параметры:
\begin{itemize}
    \item $\sigma_s$ --- коэффициент рассеяния (scattering coefficient)
    \item $\sigma_a$ --- коэффициент поглощения (absorption coefficient)
    \item $\sigma_t = \sigma_s + \sigma_a$ --- коэффициент экстинкции (extinction coefficient)
    \item $g$ --- параметр анизотропии фазовой функции Хеньи-Грина
\end{itemize}

\subsection{Shader Storage Buffer Objects (SSBO) --- формат std430}

SSBO используют более компактное выравнивание std430, которое не требует padding между элементами массивов.

\subsubsection{Буфер материалов}

Каждый материал занимает 20 float32 (80 байт):

\begin{align}
\text{matData}[i] &= \begin{bmatrix}
\text{type} & \text{(float, индекс 0)} \\
\text{rough} & \text{(float, индекс 1)} \\
\text{ior} & \text{(float, индекс 2)} \\
\text{smoothness} & \text{(float, индекс 3)} \\
\text{albedo.r, albedo.g, albedo.b} & \text{(float, индексы 4-6)} \\
\text{reflectivity} & \text{(float, индекс 7)} \\
\text{emit.r, emit.g, emit.b} & \text{(float, индексы 8-10)} \\
\text{pad0} & \text{(float, индекс 11)} \\
\text{absorption.r, absorption.g, absorption.b} & \text{(float, индексы 12-14)} \\
\text{absorptionScale} & \text{(float, индекс 15)} \\
\text{tint.r, tint.g, tint.b} & \text{(float, индексы 16-18)} \\
\text{pad2} & \text{(float, индекс 19)}
\end{bmatrix}
\end{align}

Размер буфера материалов:
\[
\text{size}_{\text{materials}} = N_{\text{materials}} \times 20 \times 4 \text{ байт} = 80N_{\text{materials}} \text{ байт}
\]

где $N_{\text{materials}}$ --- количество материалов в сцене.

\subsubsection{Буфер объектов}

Каждый объект занимает 12 float32 (48 байт):

\begin{align}
\text{objData}[i] &= \begin{bmatrix}
\text{type} & \text{(float, индекс 0)} \\
\text{materialIndex} & \text{(float, индекс 1)} \\
\text{pad0, pad1} & \text{(float, индексы 2-3)} \\
\text{pos.x, pos.y, pos.z} & \text{(float, индексы 4-6)} \\
\text{pad2} & \text{(float, индекс 7)} \\
\text{size.x, size.y, size.z} & \text{(float, индексы 8-10)} \\
\text{pad3} & \text{(float, индекс 11)}
\end{bmatrix}
\end{align}

Размер буфера объектов:
\[
\text{size}_{\text{objects}} = N_{\text{objects}} \times 12 \times 4 \text{ байт} = 48N_{\text{objects}} \text{ байт}
\]

\subsubsection{Буфер индексов источников света}

Список индексов эмиссивных объектов хранится как массив int32:

\[
\text{lightIndices} = [i_1, i_2, \ldots, i_{N_{\text{lights}}}]
\]

где $i_j$ --- индекс объекта-источника света в массиве объектов.

Размер:
\[
\text{size}_{\text{lights}} = N_{\text{lights}} \times 4 \text{ байт}
\]

\subsubsection{Накопительный буфер}

Накопительный буфер хранит сумму всех сэмплов для каждого пикселя:

\[
\text{accumData}[i] = [R_i, G_i, B_i]
\]

где $i = y \times W + x$ --- линейный индекс пикселя с координатами $(x, y)$.

Размер накопительного буфера:
\[
\text{size}_{\text{accum}} = W \times H \times 3 \times 4 \text{ байт} = 12WH \text{ байт}
\]

где $W$ --- ширина изображения, $H$ --- высота изображения.

\section{Распараллеливание вычислений}

\subsection{Work Group размер}

Compute shader использует фиксированный размер рабочей группы (work group):

\[
\text{local\_size\_x} = 16, \quad \text{local\_size\_y} = 16
\]

Это означает, что каждая рабочая группа содержит $16 \times 16 = 256$ потоков, которые выполняются параллельно на одном вычислительном блоке GPU.

\subsection{Диспетчеризация compute shader}

Количество рабочих групп вычисляется по формулам:

\begin{align}
\text{groupsX} &= \left\lceil \frac{W}{16} \right\rceil = \left\lfloor \frac{W + 15}{16} \right\rfloor \\
\text{groupsY} &= \left\lceil \frac{H}{16} \right\rceil = \left\lfloor \frac{H + 15}{16} \right\rfloor
\end{align}

где $\lceil \cdot \rceil$ --- функция округления вверх (ceiling), $\lfloor \cdot \rfloor$ --- округление вниз (floor).

Общее количество рабочих групп:
\[
N_{\text{groups}} = \text{groupsX} \times \text{groupsY}
\]

Общее количество потоков (может превышать количество пикселей):
\[
N_{\text{threads}} = N_{\text{groups}} \times 256
\]

\subsection{Глобальный индекс пикселя}

Каждый поток имеет уникальный глобальный индекс, вычисляемый как:

\begin{align}
\text{pix}_x &= \text{gl\_GlobalInvocationID.x} \\
\text{pix}_y &= \text{gl\_GlobalInvocationID.y}
\end{align}

где $\text{gl\_GlobalInvocationID}$ --- встроенная переменная OpenGL, содержащая глобальные координаты потока.

Линейный индекс пикселя:
\[
\text{pixIdx} = \text{pix}_y \times W + \text{pix}_x
\]

\subsection{Инициализация генератора случайных чисел}

Для обеспечения детерминированности и воспроизводимости результатов, каждый поток инициализирует свой генератор случайных чисел на основе координат пикселя и seed кадра:

\[
\text{state}_0 = \text{hash}(\text{pix}_x \times 1973 \oplus \text{pix}_y \times 9277 \oplus \text{frameSeed})
\]

где $\oplus$ --- операция XOR, $\text{hash}$ --- хэш-функция для генерации псевдослучайных чисел.

Хэш-функция реализуется как последовательность битовых операций:

\begin{align}
x_1 &= x \oplus (x \gg 17) \\
x_2 &= x_1 \times 0\text{xED5AD4BB} \\
x_3 &= x_2 \oplus (x_2 \gg 11) \\
x_4 &= x_3 \times 0\text{xAC4C1B51} \\
x_5 &= x_4 \oplus (x_4 \gg 15) \\
\text{hash}(x) &= (x_5 \times 0\text{x31848BAB}) \oplus (x_5 \gg 14)
\end{align}

где $\gg$ --- побитовый сдвиг вправо.

\section{Генерация первичных лучей}

\subsection{Построение системы координат камеры}

Система координат камеры определяется тремя ортогональными векторами:

\begin{align}
\mathbf{w} &= \frac{\mathbf{c}_{\text{pos}} - \mathbf{c}_{\text{target}}}{\|\mathbf{c}_{\text{pos}} - \mathbf{c}_{\text{target}}\|} \\
\mathbf{u} &= \frac{\mathbf{c}_{\text{up}} \times \mathbf{w}}{\|\mathbf{c}_{\text{up}} \times \mathbf{w}\|} \\
\mathbf{v} &= \mathbf{w} \times \mathbf{u}
\end{align}

где $\mathbf{w}$ направлен от цели к позиции камеры (обратное направление взгляда), $\mathbf{u}$ --- вектор вправо, $\mathbf{v}$ --- вектор вверх.

\subsection{Вычисление размеров viewport}

Размеры viewport вычисляются на основе поля зрения (FOV) и соотношения сторон:

\begin{align}
\theta &= \frac{\text{FOV} \times \pi}{180} \quad \text{(преобразование из градусов в радианы)} \\
h &= \tan\left(\frac{\theta}{2}\right) \\
\text{viewportHeight} &= 2h \\
\text{viewportWidth} &= \text{aspect} \times \text{viewportHeight}
\end{align}

где $\text{aspect} = W/H$ --- соотношение сторон изображения.

\subsection{Построение луча без глубины резкости}

Для простой камеры (без DOF) луч строится следующим образом:

\begin{align}
u &= \frac{\text{pix}_x + \xi_x}{W - 1} \quad \text{(нормализованные координаты [0, 1])} \\
v &= \frac{H - 1 - \text{pix}_y + \xi_y}{H - 1}
\end{align}

где $\xi_x, \xi_y \in [0, 1]$ --- случайные смещения для сэмплирования внутри пикселя.

Направление луча:
\begin{align}
\mathbf{horizontal} &= \text{viewportWidth} \times \text{focusDist} \times \mathbf{u} \\
\mathbf{vertical} &= \text{viewportHeight} \times \text{focusDist} \times \mathbf{v} \\
\mathbf{lowerLeftCorner} &= \mathbf{c}_{\text{pos}} - \frac{1}{2}\mathbf{horizontal} - \frac{1}{2}\mathbf{vertical} - \text{focusDist} \times \mathbf{w} \\
\mathbf{dir} &= \mathbf{lowerLeftCorner} + u \times \mathbf{horizontal} + v \times \mathbf{vertical} - \mathbf{c}_{\text{pos}} \\
\mathbf{r}_{\text{dir}} &= \frac{\mathbf{dir}}{\|\mathbf{dir}\|}
\end{align}

Начало луча:
\[
\mathbf{r}_{\text{orig}} = \mathbf{c}_{\text{pos}}
\]

\subsection{Генерация луча с учетом глубины резкости (DOF)}

Для камеры с диафрагмой (aperture $> 0$) луч генерируется с учетом глубины резкости:

\begin{align}
\text{lensRadius} &= \frac{\text{aperture}}{2} \\
\mathbf{rd} &= \text{lensRadius} \times \text{randomInUnitSphere}() \\
\mathbf{offset} &= \mathbf{u} \times \mathbf{rd}_x + \mathbf{v} \times \mathbf{rd}_y \\
\mathbf{r}_{\text{orig}} &= \mathbf{c}_{\text{pos}} + \mathbf{offset} \\
\mathbf{dir} &= \mathbf{lowerLeftCorner} + u \times \mathbf{horizontal} + v \times \mathbf{vertical} - \mathbf{r}_{\text{orig}} \\
\mathbf{r}_{\text{dir}} &= \frac{\mathbf{dir}}{\|\mathbf{dir}\|}
\end{align}

\subsection{Stratified Sampling}

Для уменьшения дисперсии (variance) используется стратифицированное сэмплирование. Пиксель разбивается на $N \times N$ страт (обычно $N = 4$, что дает 16 страт):

\begin{align}
\text{strataSize} &= 4 \\
\text{totalStrata} &= \text{strataSize}^2 = 16 \\
\text{samplesPerStratum} &= \left\lfloor \frac{N_{\text{samples}}}{\text{totalStrata}} \right\rfloor \\
\text{extraSamples} &= N_{\text{samples}} - \text{samplesPerStratum} \times \text{totalStrata}
\end{align}

Для каждой страты $(s_x, s_y)$ генерируются случайные смещения:

\begin{align}
\xi_x, \xi_y &\sim \mathcal{U}(0, 1) \quad \text{(равномерное распределение)} \\
\text{stratumU} &= \frac{s_x + \xi_x}{\text{strataSize}} \\
\text{stratumV} &= \frac{s_y + \xi_y}{\text{strataSize}} \\
u &= \frac{\text{pix}_x + \text{stratumU}}{W - 1} \\
v &= \frac{H - 1 - \text{pix}_y + \text{stratumV}}{H - 1}
\end{align}

\section{Пересечение лучей с объектами}

\subsection{Пересечение со сферой}

Луч задается параметрическим уравнением:
\[
\mathbf{r}(t) = \mathbf{o} + t\mathbf{d}, \quad t > 0
\]

где $\mathbf{o}$ --- начало луча, $\mathbf{d}$ --- направление луча (нормализовано).

Сфера задается центром $\mathbf{c}$ и радиусом $r$. Подставляя уравнение луча в уравнение сферы:

\[
\|\mathbf{r}(t) - \mathbf{c}\|^2 = r^2
\]

Раскрывая скобки:

\begin{align}
\|\mathbf{o} + t\mathbf{d} - \mathbf{c}\|^2 &= r^2 \\
\|\mathbf{o} - \mathbf{c} + t\mathbf{d}\|^2 &= r^2 \\
(\mathbf{o} - \mathbf{c} + t\mathbf{d}) \cdot (\mathbf{o} - \mathbf{c} + t\mathbf{d}) &= r^2
\end{align}

Раскрывая скалярное произведение:

\begin{align}
\|\mathbf{o} - \mathbf{c}\|^2 + 2t(\mathbf{o} - \mathbf{c}) \cdot \mathbf{d} + t^2\|\mathbf{d}\|^2 &= r^2
\end{align}

Поскольку $\mathbf{d}$ нормализован, $\|\mathbf{d}\|^2 = 1$. Получаем квадратное уравнение:

\begin{align}
at^2 + bt + c &= 0
\end{align}

где:
\begin{align}
a &= \|\mathbf{d}\|^2 = 1 \\
b &= 2(\mathbf{o} - \mathbf{c}) \cdot \mathbf{d} \\
c &= \|\mathbf{o} - \mathbf{c}\|^2 - r^2
\end{align}

Дискриминант:
\[
\Delta = b^2 - 4ac = b^2 - 4c
\]

Если $\Delta < 0$, пересечений нет. Иначе:

\[
t = \frac{-b \pm \sqrt{\Delta}}{2a} = \frac{-b \pm \sqrt{\Delta}}{2}
\]

Выбирается минимальное положительное $t$ в диапазоне $[t_{\min}, t_{\max}]$.

Нормаль в точке пересечения:
\[
\mathbf{n} = \frac{\mathbf{r}(t) - \mathbf{c}}{r}
\]

\subsection{Пересечение с плоскостью}

Плоскость задается точкой $\mathbf{p}$ и нормалью $\mathbf{n}$ (нормализована). Уравнение плоскости:

\[
(\mathbf{x} - \mathbf{p}) \cdot \mathbf{n} = 0
\]

Подставляя параметрическое уравнение луча:

\begin{align}
(\mathbf{o} + t\mathbf{d} - \mathbf{p}) \cdot \mathbf{n} &= 0 \\
(\mathbf{o} - \mathbf{p}) \cdot \mathbf{n} + t(\mathbf{d} \cdot \mathbf{n}) &= 0
\end{align}

Отсюда:

\[
t = \frac{(\mathbf{p} - \mathbf{o}) \cdot \mathbf{n}}{\mathbf{d} \cdot \mathbf{n}}
\]

Если $|\mathbf{d} \cdot \mathbf{n}| < \epsilon$ (луч параллелен плоскости), пересечения нет. Иначе проверяется, что $t \in [t_{\min}, t_{\max}]$.

Точка пересечения:
\[
\mathbf{hit} = \mathbf{o} + t\mathbf{d}
\]

Нормаль: $\mathbf{n}$ (уже известна).

\subsection{Пересечение с кубом (AABB)}

Ограничивающий параллелепипед (Axis-Aligned Bounding Box, AABB) задается двумя точками $\mathbf{b}_{\min}$ и $\mathbf{b}_{\max}$.

Для каждой оси $i \in \{x, y, z\}$ вычисляются параметры пересечения:

\begin{align}
t_{\text{near},i} &= \frac{b_{\min,i} - o_i}{d_i} \\
t_{\text{far},i} &= \frac{b_{\max,i} - o_i}{d_i}
\end{align}

Если $d_i < 0$, луч идет в обратном направлении, поэтому значения меняются местами:

\begin{align}
\text{if } d_i < 0: \quad &t_{\text{near},i} \leftrightarrow t_{\text{far},i}
\end{align}

Итоговые параметры пересечения:

\begin{align}
t_{\text{near}} &= \max(t_{\text{near},x}, t_{\text{near},y}, t_{\text{near},z}) \\
t_{\text{far}} &= \min(t_{\text{far},x}, t_{\text{far},y}, t_{\text{far},z})
\end{align}

Пересечение существует, если:
\[
t_{\text{near}} < t_{\text{far}} \quad \text{и} \quad t_{\text{far}} > 0
\]

Параметр пересечения: $t = t_{\text{near}}$ (если $t_{\text{near}} > 0$) или $t = t_{\text{far}}$ (если луч начинается внутри куба).

Нормаль определяется по ближайшей грани:

\[
\mathbf{n}_i = \begin{cases}
(1, 0, 0) & \text{если ближайшая грань по } x \\
(0, 1, 0) & \text{если ближайшая грань по } y \\
(0, 0, 1) & \text{если ближайшая грань по } z
\end{cases}
\]

с соответствующим знаком в зависимости от направления.

\section{Материалы и BRDF}

\subsection{Уравнение рендеринга}

Уравнение рендеринга описывает количество света, исходящего из точки $\mathbf{x}$ в направлении $\omega_o$:

\[
L_o(\mathbf{x}, \omega_o) = L_e(\mathbf{x}, \omega_o) + \int_{\Omega} f_r(\mathbf{x}, \omega_i, \omega_o) L_i(\mathbf{x}, \omega_i) (\mathbf{n} \cdot \omega_i) d\omega_i
\]

где:
\begin{itemize}
    \item $L_o$ --- исходящая яркость (outgoing radiance)
    \item $L_e$ --- эмиссия (emission)
    \item $f_r$ --- BRDF (Bidirectional Reflectance Distribution Function)
    \item $L_i$ --- входящая яркость (incoming radiance)
    \item $\mathbf{n}$ --- нормаль поверхности
    \item $\Omega$ --- полусфера направлений над поверхностью
\end{itemize}

\subsection{Lambertian (диффузный) материал}

Lambertian BRDF моделирует идеально рассеивающую поверхность:

\[
f_r(\omega_i, \omega_o) = \frac{\rho}{\pi}
\]

где $\rho$ --- альбедо (albedo), цвет материала.

Это означает, что поверхность отражает свет равномерно во все направления, независимо от направления падения и наблюдения.

\subsection{Металлический материал --- GGX распределение}

Для металлических материалов используется микрофасетная модель с распределением нормалей GGX (Trowbridge-Reitz):

\[
D(\mathbf{h}) = \frac{\alpha^2}{\pi((\mathbf{n} \cdot \mathbf{h})^2(\alpha^2 - 1) + 1)^2}
\]

где:
\begin{itemize}
    \item $\mathbf{h} = \frac{\omega_i + \omega_o}{\|\omega_i + \omega_o\|}$ --- половинный вектор (half vector)
    \item $\alpha = \text{roughness}^2$ --- параметр шероховатости
    \item $\mathbf{n}$ --- нормаль поверхности
\end{itemize}

Для идеального зеркала (mirror) используется простое отражение:

\[
\omega_o = \text{reflect}(\omega_i, \mathbf{n}) = \omega_i - 2(\omega_i \cdot \mathbf{n})\mathbf{n}
\]

\subsection{Диэлектрический материал --- закон Снеллиуса и Френель}

Для прозрачных материалов (стекло, вода) используется закон Снеллиуса:

\[
\eta_1 \sin\theta_1 = \eta_2 \sin\theta_2
\]

где $\eta_1, \eta_2$ --- показатели преломления сред, $\theta_1, \theta_2$ --- углы падения и преломления.

Вектор преломления вычисляется как:

\begin{align}
\cos\theta_1 &= -\omega_i \cdot \mathbf{n} \\
\sin\theta_1 &= \sqrt{1 - \cos^2\theta_1} \\
\sin\theta_2 &= \frac{\eta_1}{\eta_2} \sin\theta_1
\end{align}

Если $\sin\theta_2 > 1$, происходит полное внутреннее отражение (TIR), и луч отражается.

Иначе:

\begin{align}
\cos\theta_2 &= \sqrt{1 - \sin^2\theta_2} \\
\mathbf{r}_{\perp} &= \eta(\omega_i + \cos\theta_1 \mathbf{n}) \\
\mathbf{r}_{\parallel} &= -\sqrt{1 - \|\mathbf{r}_{\perp}\|^2} \mathbf{n} \\
\omega_o &= \mathbf{r}_{\perp} + \mathbf{r}_{\parallel}
\end{align}

где $\eta = \eta_1 / \eta_2$ (при входе) или $\eta = \eta_2 / \eta_1$ (при выходе).

Коэффициент отражения Френеля (аппроксимация Шлика):

\[
R(\theta) = R_0 + (1 - R_0)(1 - \cos\theta)^5
\]

где:

\[
R_0 = \left(\frac{\eta_2 - \eta_1}{\eta_2 + \eta_1}\right)^2
\]

Вероятность отражения используется для выбора между отражением и преломлением методом Монте-Карло.

\subsection{Поглощение света (Beer-Lambert закон)}

При прохождении через диэлектрик свет поглощается по закону Бугера-Ламберта:

\[
I(d) = I_0 \exp(-\sigma d)
\]

где:
\begin{itemize}
    \item $I_0$ --- начальная интенсивность
    \item $\sigma = \boldsymbol{\alpha} \times \text{absorptionScale}$ --- коэффициент поглощения (вектор для RGB)
    \item $d$ --- пройденное расстояние
\end{itemize}

В коде это реализуется как:

\[
\text{attenuation} = \exp(-\boldsymbol{\alpha} \times \text{scale} \times d)
\]

\section{Importance Sampling}

\subsection{Cosine-weighted sampling для диффузных материалов}

Для эффективного сэмплирования диффузных материалов используется косинусное распределение:

\[
p(\omega) = \frac{\cos\theta}{\pi}
\]

где $\theta$ --- угол между направлением и нормалью.

Сэмплирование в локальной системе координат:

\begin{align}
\phi &= 2\pi \xi_1 \quad \text{(равномерно по азимуту)} \\
\cos\theta &= \sqrt{\xi_2} \quad \text{(косинусное распределение)} \\
\sin\theta &= \sqrt{1 - \xi_2}
\end{align}

где $\xi_1, \xi_2 \sim \mathcal{U}(0, 1)$.

Направление в локальной системе:

\[
\mathbf{d}_{\text{local}} = (\sin\theta \cos\phi, \sin\theta \sin\phi, \cos\theta)
\]

Преобразование в мировую систему через ортонормированный базис $(\mathbf{u}, \mathbf{v}, \mathbf{n})$:

\[
\mathbf{d}_{\text{world}} = \mathbf{d}_{\text{local},x} \mathbf{u} + \mathbf{d}_{\text{local},y} \mathbf{v} + \mathbf{d}_{\text{local},z} \mathbf{n}
\]

PDF по направлению:

\[
p(\omega) = \frac{\cos\theta}{\pi}
\]

\subsection{GGX importance sampling для металлов}

Для GGX распределения используется importance sampling нормалей микроповерхности:

\begin{align}
\cos\theta_h &= \sqrt{\frac{1 - \xi_2}{1 + (\alpha^2 - 1)\xi_2}} \\
\sin\theta_h &= \sqrt{1 - \cos^2\theta_h} \\
\phi_h &= 2\pi \xi_1
\end{align}

Половинный вектор в локальной системе:

\[
\mathbf{h}_{\text{local}} = (\sin\theta_h \cos\phi_h, \sin\theta_h \sin\phi_h, \cos\theta_h)
\]

Преобразование в мировую систему и вычисление направления отражения:

\[
\omega_o = \text{reflect}(-\omega_i, \mathbf{h})
\]

\subsection{Next Event Estimation (NEE)}

Для уменьшения дисперсии используется прямое сэмплирование источников света (Next Event Estimation).

Вклад от источника света:

\[
L_{\text{direct}} = \frac{f_r(\omega_i, \omega_o) L_e \cos\theta_i \cos\theta_o}{r^2 \cdot p_A}
\]

где:
\begin{itemize}
    \item $L_e$ --- эмиссия источника
    \item $\cos\theta_i$ --- косинус угла между нормалью поверхности и направлением к источнику
    \item $\cos\theta_o$ --- косинус угла между нормалью источника и направлением от источника
    \item $r$ --- расстояние до источника
    \item $p_A = 1/A$ --- PDF по площади источника (для равномерного сэмплирования)
\end{itemize}

Переход от PDF по площади к PDF по направлению:

\[
p(\omega) = \frac{p_A \cos\theta_o}{r^2}
\]

\subsection{Сэмплирование точки на сфере-источнике}

Для сферического источника света точка сэмплируется равномерно по поверхности:

\begin{align}
z &= 1 - 2\xi_1 \quad \text{(равномерно по } z \text{)} \\
r &= \sqrt{1 - z^2} \\
\phi &= 2\pi \xi_2 \\
\mathbf{p}_{\text{local}} &= (r\cos\phi, r\sin\phi, z) \\
\mathbf{p}_{\text{world}} &= \mathbf{c} + R \cdot \frac{\mathbf{p}_{\text{local}}}{\|\mathbf{p}_{\text{local}}\|}
\end{align}

где $R$ --- радиус сферы, $\mathbf{c}$ --- центр сферы.

PDF по площади:

\[
p_A = \frac{1}{4\pi R^2}
\]

\section{Path Tracing алгоритм}

\subsection{Рекурсивное уравнение рендеринга}

Path tracing решает уравнение рендеринга методом Монте-Карло, трассируя случайные пути света:

\[
L_o(\mathbf{x}, \omega_o) = L_e(\mathbf{x}, \omega_o) + \int_{\Omega} f_r(\mathbf{x}, \omega_i, \omega_o) L_i(\mathbf{x}, \omega_i) (\mathbf{n} \cdot \omega_i) d\omega_i
\]

Метод Монте-Карло аппроксимирует интеграл:

\[
L_o \approx L_e + \frac{1}{N} \sum_{i=1}^{N} \frac{f_r(\omega_i, \omega_o) L_i(\omega_i) (\mathbf{n} \cdot \omega_i)}{p(\omega_i)}
\]

где $p(\omega_i)$ --- PDF сэмплирования направления.

\subsection{Итеративный алгоритм path tracing}

Вместо рекурсии используется итеративный подход:

\begin{align}
\text{throughput} &= (1, 1, 1) \quad \text{(начальное значение)} \\
\text{radiance} &= (0, 0, 0) \\
\text{depth} &= \text{maxDepth}
\end{align}

Для каждого отскока:

\begin{align}
\text{radiance} &\leftarrow \text{radiance} + \text{throughput} \times L_e \quad \text{(если материал эмиссивный)} \\
\text{radiance} &\leftarrow \text{radiance} + \text{throughput} \times L_{\text{direct}} \quad \text{(NEE)} \\
\omega_i &\sim p(\omega) \quad \text{(сэмплирование нового направления)} \\
\text{throughput} &\leftarrow \text{throughput} \times \frac{f_r(\omega_i, \omega_o) (\mathbf{n} \cdot \omega_i)}{p(\omega_i)} \\
\text{depth} &\leftarrow \text{depth} - 1
\end{align}

\subsection{Russian Roulette}

Для раннего прекращения трассировки используется Russian Roulette:

\[
\text{continue} = \begin{cases}
\text{true} & \text{с вероятностью } p = \min(\max(\text{throughput}), 0.95) \\
\text{false} & \text{с вероятностью } 1 - p
\end{cases}
\]

Если трассировка продолжается, throughput корректируется:

\[
\text{throughput} \leftarrow \frac{\text{throughput}}{p}
\]

Это обеспечивает несмещенность оценки (unbiased estimation).

\subsection{Условия прекращения}

Трассировка прекращается, если:
\begin{itemize}
    \item $\text{depth} \leq 0$ --- достигнута максимальная глубина
    \item Луч не пересекает объекты (уходит в небо)
    \item Russian Roulette решает прекратить трассировку
    \item $\max(\text{throughput}) < \epsilon$ --- throughput стал слишком мал
\end{itemize}

\section{Объемное рассеяние (Volumetric Scattering)}

\subsection{Уравнение переноса излучения}

Уравнение переноса излучения описывает изменение яркости вдоль луча в объемной среде:

\[
\frac{dL(\mathbf{x}, \omega)}{ds} = -\sigma_t(\mathbf{x}) L(\mathbf{x}, \omega) + \sigma_s(\mathbf{x}) \int_{4\pi} p(\omega_i, \omega) L_i(\mathbf{x}, \omega_i) d\omega_i
\]

где:
\begin{itemize}
    \item $\sigma_t = \sigma_s + \sigma_a$ --- коэффициент экстинкции
    \item $\sigma_s$ --- коэффициент рассеяния
    \item $\sigma_a$ --- коэффициент поглощения
    \item $p(\omega_i, \omega)$ --- фазовая функция
\end{itemize}

\subsection{Фазовая функция Хеньи-Грина}

Фазовая функция Хеньи-Грина (Henyey-Greenstein) моделирует рассеяние в объемной среде:

\[
p(\cos\theta) = \frac{1 - g^2}{4\pi(1 + g^2 - 2g\cos\theta)^{3/2}}
\]

где:
\begin{itemize}
    \item $g \in [-1, 1]$ --- параметр анизотропии
    \item $g = 0$ --- изотропное рассеяние
    \item $g > 0$ --- рассеяние вперед (forward scattering)
    \item $g < 0$ --- рассеяние назад (backward scattering)
\end{itemize}

\subsection{Single Scattering интеграл}

Для single scattering (однократное рассеяние) интеграл упрощается. Вклад от источника света в точке $\mathbf{x}$:

\[
L_s(\mathbf{x}, \omega) = \int_0^d T(\mathbf{x}, \mathbf{x} + s\omega) \sigma_s(\mathbf{x} + s\omega) \int_{4\pi} p(\omega_i, \omega) L_e(\mathbf{x} + s\omega, \omega_i) d\omega_i ds
\]

где $T$ --- функция пропускания (transmittance):

\[
T(\mathbf{a}, \mathbf{b}) = \exp\left(-\int_0^{\|\mathbf{b} - \mathbf{a}\|} \sigma_t(\mathbf{a} + t\frac{\mathbf{b} - \mathbf{a}}{\|\mathbf{b} - \mathbf{a}\|}) dt\right)
\]

Для однородной среды:

\[
T(d) = \exp(-\sigma_t d)
\]

\subsection{Численное интегрирование}

Интеграл аппроксимируется методом прямоугольников:

\[
L_s \approx \sum_{i=0}^{N-1} T(s_i) \sigma_s(\mathbf{x}_i) p(\omega_i, \omega) L_e(\mathbf{x}_i, \omega_i) \Delta s
\]

где:
\begin{itemize}
    \item $s_i = (i + 0.5)\Delta s$ --- позиция сэмпла
    \item $\Delta s = d/N$ --- шаг интегрирования
    \item $\mathbf{x}_i = \mathbf{x} + s_i \omega$ --- позиция сэмпла
    \item $N$ --- количество шагов (обычно 24)
\end{itemize}

\subsection{Beer-Lambert закон для поглощения}

При прохождении через среду свет ослабляется по закону Бугера-Ламберта:

\[
L(d) = L_0 \exp(-\sigma_a d)
\]

где $d$ --- пройденное расстояние.

Для неоднородной среды:

\[
L(d) = L_0 \exp\left(-\int_0^d \sigma_a(\mathbf{x} + t\omega) dt\right)
\]

\section{Накопление и усреднение сэмплов}

\subsection{Накопление сэмплов}

Для каждого пикселя накапливается сумма всех сэмплов:

\[
\text{accum}[i] = \sum_{s=1}^{N} L_s(i)
\]

где:
\begin{itemize}
    \item $i = y \times W + x$ --- линейный индекс пикселя
    \item $L_s(i)$ --- яркость $s$-го сэмпла для пикселя $i$
    \item $N$ --- общее количество сэмплов
\end{itemize}

В коде это реализуется как:

\begin{align}
\text{accumData}[i \times 3 + 0] &\leftarrow \text{accumData}[i \times 3 + 0] + L_s(i).r \\
\text{accumData}[i \times 3 + 1] &\leftarrow \text{accumData}[i \times 3 + 1] + L_s(i).g \\
\text{accumData}[i \times 3 + 2] &\leftarrow \text{accumData}[i \times 3 + 2] + L_s(i).b
\end{align}

\subsection{Усреднение}

Итоговый цвет пикселя вычисляется как среднее арифметическое:

\[
\bar{L}(i) = \frac{1}{N} \sum_{s=1}^{N} L_s(i) = \frac{\text{accum}[i]}{N}
\]

где $N = \text{sampleCount}$ --- текущее количество накопленных сэмплов.

Компоненты цвета:

\begin{align}
\bar{L}_r(i) &= \frac{\text{accumData}[i \times 3 + 0]}{N} \\
\bar{L}_g(i) &= \frac{\text{accumData}[i \times 3 + 1]}{N} \\
\bar{L}_b(i) &= \frac{\text{accumData}[i \times 3 + 2]}{N}
\end{align}

\subsection{Запись в текстуру}

Итоговый цвет записывается в текстуру формата RGBA16F (16-битное число с плавающей точкой на канал):

\[
\text{imageStore}(\text{destTex}, \text{pix}, (\bar{L}_r, \bar{L}_g, \bar{L}_b, 1.0))
\]

Формат RGBA16F позволяет хранить значения вне диапазона [0, 1], что важно для HDR (High Dynamic Range) рендеринга.

\section{Оптимизации}

\subsection{Кэширование вычислений}

Для ускорения вычислений используются оптимизации:

\subsubsection{Быстрая нормализация}

Вместо стандартной нормализации используется оптимизированная версия:

\[
\mathbf{v}_{\text{norm}} = \mathbf{v} \times \text{inversesqrt}(\|\mathbf{v}\|^2)
\]

где $\text{inversesqrt}$ --- быстрая инструкция GPU для вычисления $1/\sqrt{x}$.

\subsubsection{Кэширование квадратов}

Часто используемые квадраты вычисляются один раз:

\[
r^2 = r \times r, \quad \alpha^2 = \alpha \times \alpha
\]

\subsubsection{Кэширование обратных значений}

Обратные значения вычисляются один раз:

\[
\text{invA} = \frac{1}{a}, \quad \text{invRadius} = \frac{1}{r}
\]

\subsection{Ранний выход из циклов}

Для оптимизации используются ранние выходы:

\begin{itemize}
    \item Если луч не пересекает объект, цикл прерывается
    \item Если throughput стал слишком мал, трассировка прекращается
    \item Если материал эмиссивный, трассировка прекращается после добавления эмиссии
\end{itemize}

\subsection{Оптимизация доступа к памяти}

\begin{itemize}
    \item Последовательный доступ к массивам объектов и материалов
    \item Кэширование часто используемых значений в регистрах
    \item Минимизация количества обращений к глобальной памяти
\end{itemize}

\subsection{Firefly Reduction}

Для уменьшения артефактов (fireflies) от очень ярких сэмплов используется ограничение яркости:

\[
L_{\text{clamped}} = \begin{cases}
L & \text{если } \text{luminance}(L) \leq L_{\max} \\
L \times \frac{L_{\max}}{\text{luminance}(L)} & \text{иначе}
\end{cases}
\]

где:

\[
\text{luminance}(L) = 0.2126 \times L_r + 0.7152 \times L_g + 0.0722 \times L_b
\]

и $L_{\max} = 500.0$ --- максимальная допустимая яркость.

\section{Заключение}

В данном документе представлено полное математическое обоснование работы патрейсера на видеокарте. Описаны:

\begin{itemize}
    \item Форматы данных и их передача на GPU через UBO и SSBO
    \item Механизм распараллеливания через compute shaders
    \item Математические формулы для генерации лучей и пересечений
    \item BRDF для различных типов материалов
    \item Алгоритмы importance sampling
    \item Path tracing с Russian Roulette
    \item Объемное рассеяние света
    \item Накопление и усреднение сэмплов
    \item Оптимизации для повышения производительности
\end{itemize}

Все формулы выведены из физических принципов и соответствуют реализации в коде.

\end{document}

